\section{Introduction}

\subsection{Motivation and Related Works}
Cooperative multi-quadrotor systems provide scalable, high-capacity solutions for payload transportation \cite{guebsi_drones_2024, tan_cooperative_2018}. While cable-suspended payloads \cite{zeng_decentralized_2025,peng_modeling_2023} are common, rigidly attaching quadrotors to a common frame allows the distributed system to be controlled as a single composite rigid body \cite{tan_cooperative_2018}. However, this geometric simplification assumes a perfectly symmetrical mass distribution \cite{mellinger2013trajectory, ye_nonlinear_2023}. In practice, uncertain payload Centers of Gravity (CoG) introduce unmodeled disturbance torques, causing severe lateral drift that may lead to a crash \cite{lee_concurrent_2025,farid_digital_2025}.

To mitigate these unmodeled disturbances, recent literature has explored structural hardware optimizations, such as canting rotors to artificially increase yaw authority \cite{tan_cooperative_2018}, alongside advanced algorithmic adaptations. Various parameter estimation and control architectures have been successfully deployed to handle irregular payloads, including classical batch and recursive least-squares algorithms \cite{mellinger2013trajectory}, Concurrent Learning (CL) frameworks \cite{lee_concurrent_2025}, and Digital Twin (DT) Model Predictive Control (MPC) architectures \cite{farid_digital_2025}. Building upon this established body of work, there is an opportunity to explore emerging machine learning techniques as alternative estimation mechanisms for real-time flight stabilization.

This research develops a Proof of Concept (PoC) for a Physics-Informed Neural Network (PINN) estimator tailored for cooperative, rigid multi-quadrotor payload transportation inside a simulation. Specifically, this work explores the feasibility of using a PINN to mitigate dynamic disturbances caused by an unmodeled, off-center static payload. To bridge the gap between PINN-based parameter estimation and multi-quadrotor control, the primary contributions of this work are:

\begin{enumerate}
    \item A PINN estimator that dynamically identifies unknown payload mass and static spatial shifts by utilizing a short-term receding data horizon;
    \item The embedding of an energy-based loss function derived from Euler-Lagrange and Euler-Poincaré equations; 
    \item The integration of the PINN estimator with a small angle multi-agent controller based on \cite{mellinger2013trajectory} to continuously feed updated offset parameters into the control allocation matrix, thereby restoring tracking stability.
\end{enumerate}